\chapter{Twierdzenie Taylora}
\label{ma2:ch:taylor}

W całym rozdziale $E,F$ są rzeczywistymi przestrzeniami
Banacha, $U\subseteq E$ jest otwarty i $f:U\to F$.
Korzystamy z pochodnych Frécheta i ich wieloliniowej reprezentacji
z rozdziału \ref{ma2:ch:higher-derivatives}.
Wzór całkowy rozwija wynik z AF 15
(\atlasref{FA-15::thm:taylor15}{wzór Taylora z resztą całkową}).

\section{Wielomian Taylora i reszta Peana}
\begin{definicja}
\label{ma2:def:taylor-polynomial}
Jeżeli istnieją pochodne $D^jf(p)$ dla $1\le j\le m$,
to wielomianem Taylora stopnia co najwyżej $m$ w punkcie
$p\in U$ nazywamy odwzorowanie $T_{m,p}f:E\to F$:
\[
 T_{m,p}f(h)=\sum_{j=0}^m\frac{1}{j!}D^jf(p)[h^j],
 \qquad D^0f(p)[h^0]=f(p).
\]
Tutaj $D^jf(p)[h^j]=D^jf(p)[h,\ldots,h]$.
Dla $p+h\in U$ kładziemy
$R_{m,p}f(h)=f(p+h)-T_{m,p}f(h)$.
\end{definicja}
Z wniosku \ref{ma2:cor:pointwise-higher-symmetric}
pochodne są symetryczne także w przypadku punktowej
różniczkowalności wyższego rzędu. Każdy składnik o indeksie
$j\ge1$ jest ciągłym wielomianem jednorodnym stopnia $j$
(stwierdzenie \ref{ma2:prop:homogeneous-polynomial}).

\begin{twierdzenie}[Wzór Taylora z resztą Peana]
\label{ma2:thm:taylor-peano}
Niech $m\ge1$. Jeżeli istnieje $D^mf(p)$ w sensie
iterowanej pochodnej Frécheta z definicji \ref{ma2:def:ck},
to
\[
 f(p+h)=T_{m,p}f(h)+o(\|h\|^m)\qquad(h\to0).
\]
Nie zakładamy ciągłości $D^mf$ w otoczeniu $p$ ani
istnienia $D^mf$ w innych punktach.
\end{twierdzenie}
\begin{proof}
Dla $m=1$ teza jest definicją pochodnej Frécheta.
Załóżmy, że twierdzenie obowiązuje dla rzędu $m-1$,
gdzie $m\ge2$. Z istnienia $D^mf(p)$ wynika, że
$Df$ jest określona na pewnym otwartym otoczeniu $V$
punktu $p$ i ma pochodną rzędu $m-1$ w $p$ jako
odwzorowanie o wartościach w $\mathcal L(E,F)$.
Jest to przestrzeń Banacha na mocy stwierdzenia
\ref{ma2:prop:multi-banach}. Stosujemy do $Df$
założenie indukcyjne i naturalne utożsamienie
z lematu \ref{ma2:lem:multi-identification}.
Otrzymujemy w normie operatorowej
\[
 Df(p+u)=\sum_{j=0}^{m-1}\frac{1}{j!}
 D^{j+1}f(p)[u^j,\cdot]+o(\|u\|^{m-1}).
\]
Symetria pochodnych w $p$ pozwala umieścić wolny
argument na ostatniej pozycji.
Ze stwierdzenia \ref{ma2:prop:homogeneous-polynomial}
wynika, że pochodna $h\mapsto T_{m,p}f(h)$ w $u$
jest dokładnie sumą operatorów po prawej stronie.
Dlatego funkcja $r(u)=f(p+u)-T_{m,p}f(u)$
jest różniczkowalna w otoczeniu zera, $r(0)=0$
oraz $\|Dr(u)\|=o(\|u\|^{m-1})$. Mamy również $Dr(0)=0$.

Dla dowolnego $\varepsilon>0$ wybieramy $\rho>0$ tak,
aby $p+B(0,\rho)\subseteq V$ i dla $\|u\|<\rho$
zachodziło $\|Dr(u)\|\le\varepsilon\|u\|^{m-1}$.
Jeżeli $0<\|h\|<\rho$, to na odcinku $[0,h]$
norma pochodnej $r$ nie przekracza
$\varepsilon\|h\|^{m-1}$. Funkcja jest ciągła
na odcinku dzięki różniczkowalności w jego otoczeniu.
Z twierdzenia \ref{ma2:thm:mean-value-inequality}
otrzymujemy
\[
 \|r(h)\|=\|r(h)-r(0)\|\le\varepsilon\|h\|^m.
\]
Dowolność $\varepsilon$ daje żądaną granicę.
Indukcja jest zakończona.
\end{proof}

\begin{stwierdzenie}[Jednoznaczność wielomianu przybliżającego]
\label{ma2:prop:taylor-uniqueness}
Niech $P(h)=\sum_{j=0}^m A_j[h^j]$,
gdzie $A_0\in F$, $A_j\in\mathcal L_s^j(E;F)$
dla $j\ge1$. Jeżeli $f(p+h)-P(h)=o(\|h\|^m)$
oraz istnieje $D^mf(p)$, to $A_0=f(p)$ oraz
$A_j=D^jf(p)/j!$ dla $1\le j\le m$.
\end{stwierdzenie}
\begin{proof}
Z twierdzenia \ref{ma2:thm:taylor-peano} różnica
$P(h)-T_{m,p}f(h)$ jest $o(\|h\|^m)$.
Oznaczmy jej współczynniki przez $B_j$.
Jeżeli $E=\{0\}$, przyjmujemy standardową konwencję
reszty Peana z wartością zero w $h=0$; wszystkie
odwzorowania wieloliniowe dodatniego rzędu są zerowe.
Niech dalej $E\ne\{0\}$. Dla ustalonego $v\ne0$
podstawiamy $h=tv$, $t>0$. Mamy
$\sum_{j=0}^m t^j B_j[v^j]=o(t^m)$, ponieważ
$\|tv\|^m=t^m\|v\|^m$.
Przechodząc do granicy $t\downarrow0$, dostajemy $B_0=0$.
Jeżeli $B_j[v^j]=0$ dla $j<k$, dzielimy równość
przez $t^k$ i przechodzimy do granicy.
Wszystkie wyrazy o indeksie większym od $k$ znikają,
a prawa strona dąży do zera także dla $k=m$.
Zatem $B_k[v^k]=0$. Indukcja daje znikanie wszystkich
wartości na przekątnej dla każdego $v$ (dla $v=0$
wynika to z wieloliniowości).
Dla $j\ge1$ wielomian $h\mapsto B_j[h^j]$ jest więc
zerowy. Jego pochodna rzędu $j$ w zerze wynosi $j!B_j$
na mocy stwierdzenia \ref{ma2:prop:homogeneous-polynomial}.
Stąd $B_j=0$, co daje tezę.
\end{proof}

\section{Reszta całkowa}
\begin{twierdzenie}[Wzór Taylora na odcinku]
\label{ma2:thm:taylor-integral}
Niech $m\ge1$, $f\in C^m(U,F)$ oraz $[p,p+h]\subseteq U$.
Wtedy
\[
\begin{aligned}
 f(p+h)={}&T_{m-1,p}f(h)\\
 &+\frac{1}{(m-1)!}\int_0^1(1-t)^{m-1}
                   D^mf(p+th)[h^m]\,\mathrm dt.
\end{aligned}
\]
Całka jest całką Riemanna o wartościach w $F$.
\end{twierdzenie}
\begin{proof}
Jeżeli $h=0$, wszystkie składniki o dodatnim rzędzie
są zerowe i równość zachodzi. Niech $h\ne0$.
Zbiór $I=\{t\in\R:p+th\in U\}$ jest otwarty i zawiera
$[0,1]$. Zgodnie ze stwierdzeniem \ref{ma2:prop:line-higher}
krzywa $\gamma(t)=f(p+th)$ należy do $C^m(I,F)$ oraz
$\gamma^{(j)}(t)=D^jf(p+th)[h^j]$ dla $1\le j\le m$.
Każda z całkowanych funkcji jest ciągła na $[0,1]$;
jej całka istnieje dzięki zupełności $F$.

Dowodzimy wzoru dla $\gamma$ przez indukcję po $m$.
Dla $m=1$ lemat \ref{ma2:lem:ftc-banach} daje
$\gamma(1)=\gamma(0)+\int_0^1\gamma'(t)\,\mathrm dt$.
Przyjmijmy wzór rzędu $m$ i niech $\gamma\in C^{m+1}(I,F)$.
Funkcja $t\mapsto (1-t)^m\gamma^{(m)}(t)$ jest klasy
$C^1$ na $[0,1]$; reguła iloczynu skalaru i wektora daje
jej pochodną
\[
 -m(1-t)^{m-1}\gamma^{(m)}(t)
       +(1-t)^m\gamma^{(m+1)}(t).
\]
Stosując lemat \ref{ma2:lem:ftc-banach} i uwzględniając
wartości na końcach, dostajemy
\[
\begin{aligned}
 \frac1{(m-1)!}\int_0^1(1-t)^{m-1}\gamma^{(m)}(t)\,\mathrm dt
 ={}&\frac{\gamma^{(m)}(0)}{m!}\\
 &+\frac1{m!}\int_0^1(1-t)^m\gamma^{(m+1)}(t)\,\mathrm dt .
\end{aligned}
\]
Podstawiamy tę równość do reszty we wzorze indukcyjnym.
Uzyskany wzór jest wzorem rzędu $m+1$.
Podstawienie pochodnych krzywej daje twierdzenie.
\end{proof}

\begin{wniosek}[Oszacowania obu reszt]
\label{ma2:cor:taylor-integral-bounds}
Przy założeniach twierdzenia \ref{ma2:thm:taylor-integral}
zachodzą nierówności
\[
\begin{aligned}
 \|R_{m-1,p}f(h)\|&\le
 \frac{\|h\|^m}{m!}\sup_{0\le t\le1}\|D^mf(p+th)\|,\\
 \|R_{m,p}f(h)\|&\le
 \frac{\|h\|^m}{m!}\sup_{0\le t\le1}
               \|D^mf(p+th)-D^mf(p)\|.
\end{aligned}
\]
Jeżeli ponadto na tym odcinku
$\|D^mf(p+th)-D^mf(p)\|\le L\|th\|^\beta$,
gdzie $0<\beta\le1$ i $L\ge0$, to
$\|R_{m,p}f(h)\|\le L\|h\|^{m+\beta}/m!$.
Dla $\beta=1$ można zastąpić stałą $L/m!$ przez $L/(m+1)!$.
\end{wniosek}
\begin{proof}
Pierwsze oszacowanie wynika z nierówności normy całki
przez całkę normy, uzasadnionej w dowodzie twierdzenia
\ref{ma2:thm:integral-increment}, oraz z oszacowania
wieloliniowego ze stwierdzenia \ref{ma2:prop:multi-banach}.
Waga całkowa ma całkę
$\int_0^1(1-t)^{m-1}\,\mathrm dt=1/m$.
Aby otrzymać drugie oszacowanie, odejmujemy
$D^mf(p)[h^m]/m!$ od reszty rzędu $m-1$ i wnosimy
ten składnik pod całkę:
\[
 R_{m,p}f(h)=\frac{1}{(m-1)!}\int_0^1(1-t)^{m-1}
 (D^mf(p+th)-D^mf(p))[h^m]\,\mathrm dt .
\]
Stosujemy tę samą nierówność.
Przy założeniu Höldera norma ostatniej całki nie przekracza
\[
 \frac{L\|h\|^{m+\beta}}{(m-1)!}
       \int_0^1(1-t)^{m-1}t^\beta\,\mathrm dt .
\]
Ponieważ $t^\beta\le1$ na $[0,1]$, otrzymujemy
pierwszą stałą. Dla $\beta=1$ podstawiamy $s=1-t$:
$\int_0^1 (1-t)^{m-1}t\,\mathrm dt
=1/m-1/(m+1)=1/(m(m+1))$, co daje drugą stałą.
Ciągłość $D^mf$ w $p$ oznacza, że supremum
w drugim oszacowaniu dąży do zera wraz z $h$;
każdy punkt odcinka ma odległość od $p$ co najwyżej
$\|h\|$. Otrzymujemy ponownie resztę Peana w przypadku
$f\in C^m$ bez odwoływania się do zwartości kul.
\end{proof}

\section{Reszta Lagrange'a}
\begin{twierdzenie}[Skalarna reszta Lagrange'a]
\label{ma2:thm:taylor-lagrange}
Niech $f\in C^m(U,\R)$, $m\ge1$, $h\ne0$ i
$[p,p+h]\subseteq U$. Istnieje $\theta\in(0,1)$ takie, że
\[
 f(p+h)=T_{m-1,p}f(h)+\frac{1}{m!}
                         D^mf(p+\theta h)[h^m].
\]
\end{twierdzenie}
\begin{proof}
Niech $\gamma(t)=f(p+th)$ i
$c=\gamma(1)-\sum_{j=0}^{m-1}\gamma^{(j)}(0)/j!$.
Definiujemy skalarną funkcję
\[
 \varphi(t)=\gamma(t)-\sum_{j=0}^{m-1}
             \frac{\gamma^{(j)}(0)}{j!}t^j-ct^m.
\]
Mamy $\varphi(0)=\varphi(1)=0$ oraz
$\varphi^{(j)}(0)=0$ dla $1\le j\le m-1$.
Z twierdzenia Rolle'a istnieje $\theta_1\in(0,1)$
z $\varphi'(\theta_1)=0$.
Jeżeli $m>1$, stosujemy je ponownie do $\varphi'$
na $[0,\theta_1]$; funkcja jest ciągła na przedziale,
różniczkowalna we wnętrzu i zeruje się na obu końcach.
Powtarzając procedurę dla kolejnych pochodnych, dostajemy
$\theta=\theta_m\in(0,1)$ z $\varphi^{(m)}(\theta)=0$.
Stąd $\gamma^{(m)}(\theta)=m!c$. Wzór na pochodne
krzywej ze stwierdzenia \ref{ma2:prop:line-higher}
daje żądaną równość.
\end{proof}
\begin{uwaga}
Dla wartości w przestrzeni Banacha można zastosować
poprzednie twierdzenie do $\ell\circ f$, gdzie
$\ell\in\mathcal L(F,\R)$; punkt pośredni może jednak
zależeć od $\ell$. Już dla $m=1$ przykład
\ref{ma2:ex:vector-mvt-failure} wyklucza ogólną
wektorową równość Lagrange'a.
\end{uwaga}
\begin{przyklad}[Brak wspólnego punktu pośredniego także dla $m=2$]
\label{ma2:ex:taylor-vector-lagrange}
Niech $\gamma(t)=(t^3,t^4)$ dla $t\in\R$.
Mamy $\gamma(0)=\gamma'(0)=(0,0)$ i
$\gamma(1)=(1,1)$, podczas gdy $\gamma''(t)=(6t,12t^2)$.
Równość $\gamma(1)=\gamma(0)+\gamma'(0)+
\gamma''(\theta)/2$ wymagałaby jednocześnie
$3\theta=1$ i $6\theta^2=1$. Z pierwszej równości
wynika $\theta=1/3$, lecz wtedy $6\theta^2=2/3$.
Taki punkt nie istnieje. Wzór całkowy jest natomiast
prawdziwy: całki $\int_0^1(1-t)6t\,\mathrm dt$
i $\int_0^1(1-t)12t^2\,\mathrm dt$ są obie równe $1$.
\end{przyklad}

\section{Wzór wielowskaźnikowy w skończonym wymiarze}
\begin{wniosek}
\label{ma2:cor:taylor-multiindex}
Niech $U\subseteq\R^n$ będzie otwarty,
$f\in C^m(U,F)$ oraz $[p,p+h]\subseteq U$.
Przy oznaczeniach wielowskaźnikowych z definicji
\ref{ma2:def:multiindex} mamy
\[
 f(p+h)=\sum_{|\alpha|\le m}
       \frac{\partial^\alpha f(p)}{\alpha!}h^\alpha
       +o(\|h\|^m)\qquad(h\to0).
\]
Dokładna reszta po wyrazach stopnia co najwyżej $m-1$
wynosi
\[
 m\sum_{|\alpha|=m}\frac{h^\alpha}{\alpha!}
        \int_0^1(1-t)^{m-1}\partial^\alpha f(p+th)\,\mathrm dt .
\]
\end{wniosek}
\begin{proof}
Dla $j\ge1$ wniosek \ref{ma2:cor:multiindex-diagonal}
daje $D^jf(p)[h^j]/j!
=\sum_{|\alpha|=j} h^\alpha\partial^\alpha f(p)/\alpha!$.
Dla $j=0$ przyjmujemy $\partial^0 f=f$, $h^0=1$
i $0!=1$. Sumując, dostajemy wskazany wielomian;
twierdzenie \ref{ma2:thm:taylor-peano} daje resztę Peana.
W twierdzeniu \ref{ma2:thm:taylor-integral} stosujemy
ten sam wzór do pochodnej w punkcie $p+th$.
Suma po wielowskaźnikach jest skończona, więc liniowość
całki pozwala wynieść ją przed całkę. Stosunek
$m!/(m-1)!=m$ daje wskazany współczynnik.
\end{proof}
\begin{przyklad}[Rozwinięcie funkcji dwóch zmiennych]
\label{ma2:ex:taylor-exp-product}
Niech $f(x,y)=\exp(xy)$. W punkcie $(0,0)$ mamy
$f=1$, $\partial_x f=\partial_y f=0$,
$\partial_{xx}f=\partial_{yy}f=0$ oraz
$\partial_{xy}f=1$, co wynika z
$\partial_x f=y\exp(xy)$ i $\partial_y f=x\exp(xy)$.
Wielomian Taylora rzędu drugiego to $1+xy$.
Reszta jest nawet $O(\|(x,y)\|_2^4)$: dla $s=xy$
jednowymiarowy wzór całkowy daje
$|\exp(s)-1-s|\le e^{|s|}|s|^2/2$,
a $2|xy|\le x^2+y^2$.
Dla $\|(x,y)\|_2\le1$ dostajemy zatem
$|\exp(xy)-1-xy|\le e^{1/2}\|(x,y)\|_2^4/8$.
\end{przyklad}
